\chapter{Extensions: Attention and Transformers}\label{ch:transformers}

This chapter extends the decomposition and routing-game constructions to multiplicative layers and Transformers. It distinguishes the sign-sorted monotone split inherited from \textcite{hiddenMono}, an exact convex split of the full deterministic Transformer developed here, and the conditioned routing-game rules used in the experiments. Each construction has its own domain and preserves a different set of structural properties.

Two features separate a Transformer \parencite{Vaswani2017} from the pointwise-activation networks treated earlier. First, the attention products $QK^\top$ and $AV$ multiply input-dependent quantities. Second, softmax and LayerNorm couple several coordinates. Section~\ref{sec:multiplicative_splits} develops the sign-sorted product rule, which preserves monotonicity on non-negative split inputs. Section~\ref{sec:transformer_dc} instead combines polarization, the multivariate softmax pair of Theorem~\ref{thm:softmax_dc_split}, and a reciprocal-square-root pair to preserve convexity as well. Section~\ref{sec:attention_routing_game} then studies attention as a conditioned Value-routing game. Its implementation on \ac{vit}-B/16, described in Section~\ref{sec:vit_implementation}, supplies the attribution results evaluated in Chapter~\ref{ch:experiments}; the new full Transformer DC construction is an analytical result.

\section{Related Work: Interpretability in Transformers}\label{sec:related_work_transformers}

Early analyses of Transformer head roles combined attention patterns with LRP-based importance estimates and pruning experiments \parencite{Voita2019}. \textcite{Abnar2020attention} objected that single-layer weights ignore propagation through depth, and proposed \emph{attention rollout} --- multiplying the residual-corrected attention matrices across layers --- and \emph{attention flow}, a max-flow computation on the layered attention graph. Both use a routing graph: rollout treats attention matrices as transition kernels of a walk on tokens, whereas attention flow treats their entries as edge capacities. What they omit is everything outside the attention weights --- value magnitudes, the \acp{mlp}, the sign structure of the projections --- so rollout can assign high relevance to tokens whose values contribute nothing.

A second line adapts \ac{lrp} \parencite{Bach2015} to the Transformer. \textcite{Chefer2021GenericAttention} combined \ac{lrp}-style relevance with attention gradients to propagate class-specific relevance through self-attention, later generalized to co-attention and encoder--decoder architectures \parencite{Chefer2021ICCV}. \textcite{Ali2022AttnLRP} showed that naive gradient propagation through attention and layer normalization can \emph{violate conservation}; holding the attention matrix and LayerNorm denominator constant (``detached''), while retaining linear centering, gives conservative local rules. AttnLRP \parencite{Achtibat2024AttnLRP} introduced dedicated rules for the bilinear products and non-linearities, scaling to \acp{llm}. Related LRP variants provide concept-level \parencite{Achtibat2023CRP} and magnitude-aware \parencite{vukadin2024abslrp} explanations. On the vision side, DAVE \parencite{Wrobel2024} targeted patch-grid attribution artifacts in \acp{vit} \parencite{Dosovitskiy2021vit} with \emph{conditioned} operators and spatial averaging. Separately, \textcite{darcet2024vision} mitigated high-norm artifact tokens with registers. The ResLRP preprint targets attribution instability caused by cancellation between opposing residual branches. Its residual $\gamma$-rule conserves relevance and bounds absolute relevance amplification by a factor of $1+2/\gamma$ at each backward merge, assuming a nonzero residual output and $\gamma>0$ \parencite{Berend2026ResLRP}. At \ac{llm} scale, gradient approximations of activation patching \parencite{Syed2024AttributionPatching,Kramar2024AtPStar} localize circuits against interventions \parencite{Hanna2024Faithfulness}, and attribution increasingly acts as a control signal for head identification and guided decoding \parencite{Kahardipraja2025Atlas,Komorowski2026AGD}.

These approaches motivate two separate questions. The structural question asks whether the full forward map, including the attention-generating operations, admits explicit convex streams; Sections~\ref{sec:multiplicative_splits} and~\ref{sec:transformer_dc} develop the relevant product and normalization rules. The attribution question asks how relevance travels through the Value routes when the realized attention matrix is fixed. Section~\ref{sec:attention_routing_game} formulates that conditioned question as a KL-regularized routing game and recovers the conservative attention-block rule of \textcite{Ali2022AttnLRP}. This gives a precise interpretation of a routing calculation of the kind motivated by \textcite{Abnar2020attention}, with the conditioning assumptions stated explicitly.

\section{Decomposing Multiplicative Operations}\label{sec:multiplicative_splits}

Gating layers, squeeze-and-excitation blocks, and attention multiply activations entrywise or contract tensors. The sign-sorted construction of \textcite{hiddenMono} preserves the difference and monotonicity of non-negative stream pairs. Section~\ref{sec:transformer_dc} gives a convexity-preserving product identity combining with Chapter~\ref{ch:decompositions}'s multivariate softmax pair.

\paragraph{The gated product rule.}
For $a = a^{+} - a^{-}$ and $b = b^{+} - b^{-}$ with non-negative streams, as guaranteed by Section~\ref{sec:general_decomposition_def}, expand and sort the cross terms by sign:
\begin{equation}
	(a^{+} - a^{-})(b^{+} - b^{-})
	\;=\;
	\underbrace{\bigl(a^{+}b^{+} + a^{-}b^{-}\bigr)}_{\eqqcolon\, (ab)^{+}}
	\;-\;
	\underbrace{\bigl(a^{+}b^{-} + a^{-}b^{+}\bigr)}_{\eqqcolon\, (ab)^{-}}.
	\label{eq:pair_product_rule}
\end{equation}
Both terms are non-negative and non-decreasing in each incoming stream, so multiplication preserves monotone pairs, analogously to routing $\Wsig{l}{+}$ and $\Wsig{l}{-}$ separately in Section~\ref{sec:dc_splitting_relu}. Applying~\eqref{eq:pair_product_rule} entrywise splits Hadamard layers; applying it to scalar contractions splits matrix products. Convexity need not survive products. The resulting streams are finite maxima of \emph{posynomials}, polynomials with non-negative coefficients. Posynomials are closed under addition, multiplication and positive scaling, and are monotone on the non-negative orthant.

\paragraph{Maxout multiplication networks.}
We enlarge the architecture class for the splitting theorem. A rank-$k$ Maxout unit \parencite{Goodfellow2013Maxout} takes the maximum of $k$ affine forms; $\relu$ is the $k=2$ case with one form zero.

\begin{definition}[Maxout multiplication network]\label{def:maxout_multiplication_network}
A \emph{rank-$k$ Maxout multiplication network} is a finite directed acyclic graph of \emph{input neurons}, \emph{addition neurons}, and \emph{multiplication neurons}. An addition neuron $v$ carries weight vectors $w^{uv} \in \RR^k$ on its incoming arcs and computes the rank-$k$ Maxout activation $a^v = \max_{i \in [k]} \sum_{u} w^{uv}_i a^u$; a multiplication neuron carries natural exponents $w^{uv} \in \NN$ and computes the monomial $a^v = \prod_{u} (a^u)^{w^{uv}}$. The network is called \emph{monotone} if all weights are non-negative, and \emph{input-convex} if all weights are non-negative except possibly those on arcs leaving input neurons.
\end{definition}

This class includes linear and convolutional layers, biases via constant inputs, maxima and max-pooling, residual connections, Hadamard products, and matrix products. Natural exponents avoid division and fractional powers of possibly negative quantities. Definition~\ref{def:maxout_multiplication_network}'s structural conditions give the following representation.

\begin{proposition}[Representation by maxima of posynomials, \LeanResult{prop:posynomial_representation}]\label{prop:posynomial_representation}
A monotone rank-$k$ Maxout multiplication network with non-negative intermediate quantities represents the maximum of finitely many posynomials in its non-negative input variables, and is monotone on the non-negative orthant. This representation does not imply ordinary convexity: the positive-coefficient monomial $xy$ is a counterexample. Allowing signed input-incident affine maps likewise supplies no general input-convexity guarantee in the presence of multiplication.\end{proposition}

\begin{proof}[Proof sketch]
Induct in topological order on the non-negative orthant. Coordinate projections are posynomials; non-negative scaling, addition and maxima preserve finite maxima of posynomials. Products of non-negative factors distribute over finite maxima, covering every network operation. Convexity fails for $F(x,y)=xy$: $F(1,1)=1$ exceeds $\tfrac12F(2,0)+\tfrac12F(0,2)=0$.
\end{proof}

\paragraph{The splitting theorem.}
Two identities extend Section~\ref{sec:dc_splitting_relu}'s split to this class. Maxout uses the maximum-shift identity
\begin{equation}
	\max_{j \in [k]}\bigl(a_j^{+} - a_j^{-}\bigr)
	\;=\;
	\max_{j \in [k]}\Bigl\{a_j^{+} + \sum_{i \neq j} a_i^{-}\Bigr\}
	\;-\;
	\sum_{i \in [k]} a_i^{-},
	\label{eq:maxout_split_identity}
\end{equation}
which generalizes both the $\relu$ splitting identity $\relu(a - b) = \max\{a, b\} - b$ of Section~\ref{sec:dc_splitting_relu} and the gated max-pooling rule of Section~\ref{sec:cnn_decompositions}; multiplication neurons are split by the multinomial generalization of \eqref{eq:pair_product_rule}, expanding $\prod_u (a^{u,+} - a^{u,-})^{w^{uv}}$ and sorting the resulting monomials by the parity of the number of $(-)$-stream factors.

\begin{theoremenv}[Monotone split of multiplicative \ac{relu} networks, \LeanResult{thm:multiplicative_split}]\label{thm:multiplicative_split}
Let $N$ be a (residual) \ac{relu} network containing linear, convolutional, maximum, matrix-product, and Hadamard-product components. Then, on non-negative split-input variables, $N$ has two \emph{monotone} rank-$2$ Maxout multiplication networks $N^{+}$ and $N^{-}$ satisfying
\begin{equation*}
	N(x) = N^{+}(x^+,x^-) - N^{-}(x^+,x^-), \qquad x=x^+-x^-,\quad x^\pm\ge0,
\end{equation*}
Each component is monotone as a function of its non-negative split inputs and is representable there by a finite maximum of posynomials. This gives a difference-of-monotone representation; it does not establish a difference-of-convex representation in the original inputs. If the number of expanded monomials per multiplication neuron is uniformly bounded, the sizes of the split networks are linear in the size of $N$.\end{theoremenv}

\begin{proof}[Proof sketch]
In topological order, replace each neuron by a function-preserving pair $(v^{+}, v^{-})$, extending \textcite{hertrich2025neuralnetworksvirtualextended}'s extended-formulation argument. Supply \emph{inputs} as non-negative $x^+,x^-$ (e.g.\ $\relu(x),\relu(-x)$), flipping output-weight signs toward the $(-)$-copy; unsplit inputs do not generally restore convexity with multiplication. \emph{Addition (Maxout) neurons} split weights as $\posneg{w^{uv}}{+} - \posneg{w^{uv}}{-}$ and apply~\eqref{eq:maxout_split_identity}: the maximum defines $v^{+}$ and the subtracted sum $v^{-}$, both with non-negative weights. For \emph{multiplication neurons}, expand by signs and feed one auxiliary multiplication neuron per summand into collecting addition neurons $v^{+}$ (even parity) and $v^{-}$ (odd parity). These identities give $N = N^{+} - N^{-}$; non-negative weights imply monotonicity and Proposition~\ref{prop:posynomial_representation}'s representation. Growth is controlled by the expanded monomial count: a uniform bound suffices, whereas bounded fan-in with unbounded exponents does not.
\end{proof}

Theorem~\ref{thm:multiplicative_split} preserves the computation through non-negative monotone streams. Section~\ref{sec:transformer_dc} replaces the sign-sorted product pair to preserve convexity too.

\begin{remark}[Smooth normalizers: softmax and sigmoid gates]\label{rem:smooth_gates}
Theorem~\ref{thm:multiplicative_split} covers projections, bilinear products and residual additions, but excludes softmax. Section~\ref{sec:closed_form_dc_activations} gives exact scalar sigmoid and GELU pairs; Corollary~\ref{cor:local_curvature_split} covers exponentiation exactly on a fixed bounded score interval. Splitting softmax's coupled denominator requires Theorem~\ref{thm:softmax_dc_split}'s multivariate row pair, combined next with convex products and LayerNorm.

Conditioned attribution instead fixes the realized softmax matrix $A$, making $X\mapsto AV(X)$ affine \parencite{Ali2022AttnLRP,Wrobel2024}. The routing-game sections study this conditioned calculation; the full forward decomposition retains $A$'s input dependence.
\end{remark}

\section{A Constructive DC Decomposition of Transformers}\label{sec:transformer_dc}

Replacing~\eqref{eq:pair_product_rule}'s sign-sorted product by a square identity preserves convexity and coordinatewise monotonicity on non-negative split inputs. With Theorem~\ref{thm:softmax_dc_split}'s softmax pair and a reciprocal-square-root extension, this yields a constructive DC decomposition of the full deterministic Transformer. This analytical extension is separate from the conditioned attribution implementation evaluated below.

\begin{proposition}[Convex product pair by polarization, \LeanResult{prop:convex_product_pair}]\label{prop:convex_product_pair}
For $A,B,C,D\ge0$, let
\begin{equation}
\begin{aligned}
    P(A,B,C,D)&=\tfrac12\bigl((A+C)^2+(B+D)^2\bigr),\\
    Q(A,B,C,D)&=\tfrac12\bigl((A+D)^2+(B+C)^2\bigr).
\end{aligned}
    \label{eq:polarized_product_pair}
\end{equation}
Then $P-Q=(A-B)(C-D)$. Both functions are convex, non-negative, and coordinatewise non-decreasing on $\RR^4_{\ge0}$. Hence substituting non-negative convex, coordinatewise non-decreasing incoming streams preserves these properties. In particular,
\begin{equation}
    (A-B)^2=2(A^2+B^2)-(A+B)^2
    \label{eq:convex_square_pair}
\end{equation}
is a valid convex, monotone square pair on $\RR^2_{\ge0}$.
\end{proposition}

\begin{proof}
Expansion proves recovery. Each summand squares an affine form, hence is convex and non-negative. Its partial derivatives are non-negative on the stated orthant, permitting convex non-decreasing composition with convex components. Setting $C=A$ and $D=B$ gives~\eqref{eq:convex_square_pair}.
\end{proof}

Monotonicity requires the orthant: the derivative of $(A+C)^2$ in $A$ is $2(A+C)$, possibly negative outside it. The Transformer theorem therefore uses non-negative doubled inputs, although the softmax pair has a larger domain. Apply~\eqref{eq:polarized_product_pair} to scalar products and sum corresponding components to split matrix products.

\subsection{Layer normalization with positive epsilon}

For $\varepsilon>0$ and $d\ge1$, LayerNorm on $u\in\RR^d$ uses
\begin{equation}
\begin{gathered}
    \bar u=\frac1d\sum_{j=1}^d u_j,\qquad c_i=u_i-\bar u,
    \qquad v=\frac1d\sum_{j=1}^d c_j^2\ge0,\\
    \operatorname{LN}_{\varepsilon}(u)_i
    =\gamma_i c_i(\varepsilon+v)^{-1/2}+\beta_i.
\end{gathered}
    \label{eq:layernorm_dc_definition}
\end{equation}
Centering and the learned scale and bias are signed affine; variance is a positive sum of squares. Only the reciprocal square root on non-negative arguments is new. Extending it to negative arguments gives a global convex pair when the variance streams vary independently.

\begin{proposition}[Reciprocal-square-root pair, \LeanResult{prop:reciprocal_sqrt_pair}]\label{prop:reciprocal_sqrt_pair}
Let $\varepsilon>0$, put $M=(2\varepsilon^{3/2})^{-1}$, and define
\begin{equation}
    r_{\varepsilon}(t)=
    \begin{cases}
        \varepsilon^{-1/2}-Mt,&t\le0,\\
        (\varepsilon+t)^{-1/2},&t\ge0.
    \end{cases}
    \label{eq:reciprocal_sqrt_extension}
\end{equation}
Then
\begin{equation}
    R_+(x,y)=r_{\varepsilon}(x-y)+Mx,
    \qquad R_-(x,y)=Mx
    \label{eq:reciprocal_sqrt_pair}
\end{equation}
are convex and coordinatewise non-decreasing on $\RR^2$, are non-negative on $\RR^2_{\ge0}$, and have difference $r_{\varepsilon}(x-y)$. Whenever $x-y\ge0$, this difference is $(\varepsilon+x-y)^{-1/2}$.
\end{proposition}

\begin{proof}
The extension is continuously differentiable at zero, with matching derivative $-M$. On $t>0$ its second derivative is $3/(4(\varepsilon+t)^{5/2})>0$, whereas it is affine on $t<0$. Thus $r_\varepsilon$ is convex and $-M\le r_\varepsilon'\le0$. Affine precomposition and addition preserve convexity, and
$\partial_xR_+=r_\varepsilon'+M\ge0$,
$\partial_yR_+=-r_\varepsilon'\ge0$.
The linear component $R_-$ is also monotone. Their origin values are $\varepsilon^{-1/2}$ and $0$, giving non-negativity on the orthant; cancellation proves the stated differences. Equivalently, this is the scalar lift of Proposition~\ref{prop:dc_profile_lift} with profiles $p(t)=r_\varepsilon(t)+Mt$, $q(t)=Mt$, and slope bound $M$.
\end{proof}

To decompose~\eqref{eq:layernorm_dc_definition}, first split the affine centering map. Apply~\eqref{eq:convex_square_pair} to each centered coordinate and average the two streams separately to obtain $v=v^+-v^-$. Apply~\eqref{eq:reciprocal_sqrt_pair} to $(v^+,v^-)$. Because the recovered variance is non-negative, the extension agrees exactly with the actual normalization factor. Finally use the polarized product pair to multiply each centered coordinate by that factor, followed by the signed affine scale and bias. Every intermediate stream is convex, monotone, and non-negative on the original split-input orthant. The hypothesis $\varepsilon>0$ ensures that the normalization is defined for every input, including constant vectors.

\subsection{An explicit pre-normalization architecture}

To make the arbitrary-depth claim concrete, we specify the ordinary Transformer forward map before defining its split evaluation. Tokens are indexed by $t\in\{1,\dots,T\}$, with feature vectors in $\RR^d$, where $T,d\ge1$. Every matrix and bias below is a fixed real parameter and may have either sign.

\begin{definition}[Finite pre-normalization Transformer, \LeanResult{def:prenorm_transformer}]\label{def:prenorm_transformer}
A Transformer is specified by a positional bias $P\in\RR^{T\times d}$, a finite list of blocks $B_1,\dots,B_K$, an optional final LayerNorm, and a signed affine scalar readout. In block $k$, first normalize each token,
\begin{equation*}
    U_t=\operatorname{LN}_{k,1}(X_t).
\end{equation*}
Each head $h$ has positive query/key width $d_{q,h}$ and value width $d_{v,h}$, signed affine projections, a real score scale $\tau_h$, and a fixed nonempty allowed-key set $I_{h,t}\subseteq\{1,\dots,T\}$ for every query token. It computes
\begin{equation}
\begin{gathered}
    Q_{h,t}=W_{Q,h}U_t+b_{Q,h},\qquad
    K_{h,t}=W_{K,h}U_t+b_{K,h},\\
    V_{h,t}=W_{V,h}U_t+b_{V,h},\qquad
    S_{h,ts}=\tau_h\langle Q_{h,t},K_{h,s}\rangle,\\
    A_{h,ts}=\frac{e^{S_{h,ts}}}{\sum_{r\in I_{h,t}}e^{S_{h,tr}}}
       \quad(s\in I_{h,t}),\qquad
    O_{h,t}=\sum_{s\in I_{h,t}}A_{h,ts}V_{h,s}.
\end{gathered}
    \label{eq:explicit_transformer_attention}
\end{equation}
Masked coefficients outside $I_{h,t}$ are zero. The usual scaled dot product is the choice $\tau_h=1/\sqrt{d_{q,h}}$. Summing the per-head output projections is equivalent to concatenation followed by one output projection. The residual attention and pointwise feedforward updates are
\begin{equation}
\begin{aligned}
    Y_t&=X_t+\sum_h W_{O,h}O_{h,t}+b_O,\\
    B_k(X)_t&=Y_t+W_2\varrho_k\!\left(W_1\operatorname{LN}_{k,2}(Y_t)+b_1\right)+b_2.
\end{aligned}
    \label{eq:explicit_transformer_block}
\end{equation}
Both LayerNorms have positive epsilon and arbitrary signed learned scale and bias as in~\eqref{eq:layernorm_dc_definition}. The activation $\varrho_k$ is chosen from ReLU, exact Gaussian GELU, Swish/SiLU, sigmoid, and softplus, with the positive parameters required in Chapter~\ref{ch:decompositions}. Head widths, feedforward widths, masks, and parameters may differ between blocks, while the residual width is $d$.

For any finite depth $K\ge0$, define
\begin{equation}
\begin{gathered}
    X^{(0)}=X+P,\qquad X^{(k)}=B_k(X^{(k-1)})\quad(1\le k\le K),\\
    \widehat X^{(K)}=
    \begin{cases}\operatorname{LN}_{\mathrm{out}}(X^{(K)}),&\text{if a final norm is present},\\
    X^{(K)},&\text{otherwise},\end{cases}\\
    N(X)=b_{\mathrm{out}}+\sum_{t=1}^{T}\sum_{j=1}^{d}
           w_{tj}^{\mathrm{out}}\widehat X^{(K)}_{tj}.
\end{gathered}
    \label{eq:arbitrary_depth_transformer}
\end{equation}
The optional final norm is applied tokenwise with positive epsilon. The depth-zero case has no blocks and consists of the positional affine map, the optional final norm, and the readout. Choosing zero readout weights outside one token includes a class-token logit. This is a deterministic inference computation; a fixed dropout mask can additionally be represented by fixed coordinate scalings.
\end{definition}

\subsection{The recursive split algorithm}

Expand the affine maps, dot products, and matrix contractions in Definition~\ref{def:prenorm_transformer} into scalar expressions. The expression constructors are inputs, real constants, sums, signed scalar multiples, products, supported scalar activations, softmax coordinates, and LayerNorm coordinates. Their ordinary evaluation uses the usual real functions in~\eqref{eq:explicit_transformer_attention}--\eqref{eq:arbitrary_depth_transformer}. No convexity certificate is assumed as part of an ordinary expression.

Write $\mathcal C(e)=(e^+,e^-)$ for the recursively constructed pair, as a function of independent non-negative input streams $(x,y)$. Pair addition is componentwise, and signed scalar multiplication is
\begin{equation}
    \operatorname{Scale}_c(p,q)=
    \begin{cases}(cp,cq),&c\ge0,\\(-cq,-cp),&c<0.
    \end{cases}
    \label{eq:transformer_pair_scale}
\end{equation}
The operators $\operatorname{Prod}$, $\operatorname{Softmax}_i^{\pm}$, and $\operatorname{LN}_i^{\pm}$ below are the explicit constructions of~\eqref{eq:polarized_product_pair}, Theorem~\ref{thm:softmax_dc_split}, and the LayerNorm recipe following Proposition~\ref{prop:reciprocal_sqrt_pair}. In particular, the LayerNorm pair computes the centered pairs, their square pairs, their average, the reciprocal-square-root pair, and the final product and affine parameters; it retains the input-dependent statistics.

\begin{algorithm}[htbp]
\DontPrintSemicolon
\caption{Recursive DC evaluation of a finite Transformer expression}
\label{alg:transformer_dc_compile}
\KwIn{A scalar output expression $e$ from Definition~\ref{def:prenorm_transformer}; split inputs $x,y\ge0$.}
\KwOut{A pair $\mathcal C(e)(x,y)$ whose difference is the ordinary output at $x-y$.}
For an input coordinate $i$, return $\mathcal C(X_i)=(x_i,y_i)$; for a constant $c$, return $(\max\{c,0\},\max\{-c,0\})$.\;
Recursively compute the pair of each child expression before applying its parent operation.\;
For a sum $\sum_j e_j$, return $\sum_j\mathcal C(e_j)$; for $ce_1$, return $\operatorname{Scale}_c(\mathcal C(e_1))$.\;
For $e_1e_2$, return $\operatorname{Prod}(\mathcal C(e_1),\mathcal C(e_2))$ using the four-square formula~\eqref{eq:polarized_product_pair}.\;
For a scalar activation $\varrho(e_1)$, apply its explicit pair to $\mathcal C(e_1)$.\;
For a softmax coordinate, collect the compiled allowed-score pairs and apply $\operatorname{Softmax}_i^{\pm}$ to that vector of pairs.\;
For a LayerNorm coordinate, collect the compiled feature pairs and apply $\operatorname{LN}_i^{\pm}$ with its positive epsilon and signed scale and bias.\;
Return the pair at the requested affine readout.\;
\end{algorithm}

The algorithm recurses over any finite expression in real arithmetic, independently of depth. Translation follows the finite block list: the empty list preserves its input expression; appending a block substitutes current token expressions into its formulas. Induction identifies ordinary evaluation with~\eqref{eq:arbitrary_depth_transformer}. Implementations may reuse shared expressions; correctness requires no bound on expansion or runtime.

\subsection{The full Transformer construction}

\begin{theoremenv}[Exact DC decomposition of a deterministic Transformer, \LeanResult{thm:transformer_dc}]\label{thm:transformer_dc}
Let $N\colon\RR^{T\times d}\to\RR$ be the Transformer of Definition~\ref{def:prenorm_transformer}, with any finite list of blocks, including the empty list, and identify its token-feature input with $\RR^m$, where $m=Td$. Algorithm~\ref{alg:transformer_dc_compile}, applied to its output expression, terminates and exactly decomposes this forward map. It constructs functions
$N_+,N_-\colon\RR^{2m}_{\ge0}\to\RR_{\ge0}$,
convex and coordinatewise non-decreasing on that orthant, such that
\begin{equation}
    N(x-y)=N_+(x,y)-N_-(x,y)
    \qquad(x,y\in\RR^m_{\ge0}).
    \label{eq:transformer_dc_recovery}
\end{equation}
In particular, the functions
\begin{equation}
    g(u)=N_+(\posneg{u}{+},\posneg{u}{-}),\qquad
    h(u)=N_-(\posneg{u}{+},\posneg{u}{-})
    \label{eq:transformer_raw_dc}
\end{equation}
are convex on all of $\RR^m$, are non-negative, and satisfy $N=g-h$. The functions $N_+,N_-$ are coordinatewise non-decreasing in the independent split inputs; $g,h$ need not be coordinatewise non-decreasing in the original signed variable $u$.
\end{theoremenv}

\begin{proof}
The block list yields a finite expression, so syntax recursion terminates. Structural induction shows that each compiled expression has non-negative convex, coordinatewise non-decreasing heads on the input orthant, whose difference evaluates the expression at $x-y$.

Coordinate projections and positive/negative constant splits satisfy this invariant. Sums preserve it; signed scaling rescales by a non-negative coefficient, swapping heads for negative scalars. Proposition~\ref{prop:convex_product_pair} handles products by exact recovery and convex, monotone composition on non-negative streams. The scalar pair theorems handle activations; Theorem~\ref{thm:softmax_dc_split} handles the nonempty allowed-key vector. LayerNorm follows its centering, square, reciprocal, product and affine recipe, with non-negative recovered variance and positive epsilon. Thus every constructor preserves the invariant.

Induction over the block list equates ordinary evaluation with the independently specified recurrence~\eqref{eq:arbitrary_depth_transformer}, proving~\eqref{eq:transformer_dc_recovery} at every finite depth, including zero, with an ungated signed affine readout. Finally, $u\mapsto\posneg{u}{\pm}$ has convex components valued in the orthant. Convex monotone composition gives globally convex $g,h$, and subtraction recovers $N(u)$.
\end{proof}

The construction retains input dependence through Queries, Keys, Values, softmax and LayerNorm statistics. Its real-arithmetic proof establishes a forward DC representation, without numerical-stability, attribution-quality, floating-point-error or runtime guarantees. The reciprocal pair's common coefficient $M=(2\varepsilon^{3/2})^{-1}$ grows as $\varepsilon$ shrinks; products and squares can amplify this mass, potentially requiring cancellation between enormous streams. The companion supplement's Section~\ref{app:transformer_dc_numeric} tests the same primitive pairs in a post-normalization block, whose LayerNorm placement differs from Definition~\ref{def:prenorm_transformer}. In this unstabilized float64 diagnostic, 15 of 20 cases exceed $10^{-8}$ output error and streams reach $5.95\times10^{111}$. The thesis's attribution experiments evaluate the conditioned rules below, not this convex construction.

\section{Attention via a detached-logit oracle}\label{ptnb:sec:attention}

For conditioned transport, our LayerNorm rule fixes both mean and scale, giving a pointwise affine
map. CP-LRP and DAVE instead retain linear centering while fixing the scale~\parencite{Ali2022AttnLRP,Wrobel2024}.
The QK-softmax block is replaced by an \emph{oracle game} emitting
a detached reference policy, the softmax row $A_{t\cdot}$, with occupation-measure trajectories
flowing only through the Value route. \ch{Like conservative propagation~\parencite{Ali2022AttnLRP}, this explains the network conditioned on its
realized attention pattern.}

\begin{theoremenv}[Attention, informal]\label{ptnb:thm:attention}
With reference policy $A_{t\cdot}$, the KL-regularized Value-routing subgame has an SPNE whose
occupation measure at the block input equals $\alpha\beta$-\ac{lrp} on the conditioned effective
map $x \mapsto \sum_k A_{tk}\, x_k W_v + b_v$ (Theorem~\ref{ptnb:thm:attn-ab-recovery}).
\end{theoremenv}
Our implementation applies this rule to the attention product and the Value projection as one
composite step; a full gradient fallback through the attention block degrades localization (AL
$0.624$ vs.\ \vitRGlocAL{} at the localization-selected configuration). Equivariant Reynolds
averaging optionally suppresses tokenization-grid artifacts~\parencite{Wrobel2024} (\ac{rg}+Eq). The
target rule, oracle embedding and implementation are in Section~\ref{ptnb:app:attention-subgame};
the full KL formulation and occupation-measure proof are in Chapter~\ref{app:supp_attention}
of the companion technical supplement.
\section{Extension: Vision Transformer Attention}\label{ptnb:app:vit-extension}\label{sec:attention_routing_game}\label{ptnb:app:attention-subgame}

We state the target rule, oracle interpretation and implementation used by the ViT experiments.
The full state-space formulation, KL derivation and occupation-measure proof are given in
Chapter~\ref{app:supp_attention} of the companion technical supplement; the risk-averse attention
variant is in Section~\ref{ptnb:app:risk-attn}.

\subsection{Attention as fixed-reference KL Value routing}\label{ptnb:app:attention-kl-routing}

\subsubsection{Target rule: \texorpdfstring{$\alpha\beta$-\ac{lrp}}{LRP-alpha-beta} with detached logits}\label{ptnb:app:attention-target-rule}

We first state the rule that the Routing-Game state-space construction in
Section~\ref{ptnb:app:attention-state-space} is designed to recover. We focus on a single attention head; the multi-head extension carries an additional head
index $h$ and the recovery is per-head. Note that ViT-B/16 implements the Value projection as a
single $D \times D$ linear $v = X W_v$ followed by a reshape into $n_h$ heads (the per-head matrix
$W_v^{(h)} \in \RR^{D \times d_h}$ is the head-$h$ column block of the full $W_v$), so one input
embedding $X_{k,e}$ feeds Value entries of every head and the per-head $R^X$ contributions sum at
$X_{k,e}$ to give the total. The per-head description below applies unchanged. We omit the
$\varepsilon$ stabilizer inside the attention block --- the surrounding ReLU recursion of
\S\ref{ptnb:app:routing-attribution} carries its own $\varepsilon$ on every linear-routing state and
that mechanism transfers to the $Av{+}W_v$ composite without modification, but tracking it through
the tensor indexing here only obscures the structure.

\paragraph{Forward pass.} The Value projection is applied per token: for each $k \in \{1,\dots,S\}$,
\begin{align}\label{ptnb:eq:attn-V-forward}
    v_{k,d} \;=\; \sum_e W_{v,e,d}\, X_{k,e} + b_{v,d},
    \qquad d \in \{1,\dots,d_h\},
\end{align}
with shared $W_v \in \RR^{D \times d_h}$ and $b_v \in \RR^{d_h}$ applied independently per token. Detached softmax attention $A \in \RR^{S \times S}_{\ge 0}$, with $A_{t,k} \ge 0$ and $\sum_k A_{t,k} = 1$, mixes Values across tokens independently for each feature $d$. For output token $t \in \{1,\dots,S\}$,
\begin{align}\label{ptnb:eq:attn-AV-forward}
    O_{t,d} \;=\; \sum_k A_{t,k}\, v_{k,d}
    \;=\; \sum_{k,e} A_{t,k}\, W_{v,e,d}\, X_{k,e} + b_{v,d}.
\end{align}

\paragraph{Composite linear map and $\kappa$-stream split.} Holding $A$, the LayerNorm statistics and the \acs{ptnbGELU} gates at their forward values
(\S\ref{ptnb:app:vit-implementation}) is a modeling assumption about the explanation question---given
the realized attention policy, through which Value routes does relevance flow?---and not a theorem
about relevance through $Q$ and $K$. \ch{Our attention extension is therefore a Value-only attribution
conditioned on the realized attention matrix, resembling the conservative propagation of
\textcite{Ali2022AttnLRP} (CP-LRP), which holds the attention matrix constant, and DAVE's conditioned
operators~\parencite{Wrobel2024}. It differs from AttnLRP~\parencite{Achtibat2024AttnLRP}, which also
propagates relevance through the attention-generating operations (the softmax and the bilinear QK
product) and so attributes the Query and Key paths.} With $A$ so conditioned, the attention block reduces to the
affine \emph{conditioned effective map} $X \mapsto O$ with composite weights
$W^{\mathrm{eff}}_{(k,e),\,(t,d)} = A_{t,k}\, W_{v,e,d}$ and bias $b_{v,d}$. Since $A_{t,k} \ge 0$,
the sign of the composite weight is the sign of $W_{v,e,d}$:
\begin{align}\label{ptnb:eq:attn-composite-weight}
    \bigl[\,W^{\mathrm{eff}}_{(k,e),\,(t,d)}\,\bigr]^\kappa
    \;=\; A_{t,k}\, W_{v,e,d}^\kappa,
    \qquad \kappa \in \{+,-\}.
\end{align}
For signed inputs $X$ (LayerNorm outputs), the per-token $\kappa$-stream split of the Value tensor
(as in the fused input layer of the \ac{rg}, with the bias as a bias node) and the
per-output-feature stream sum are
\begin{align}\label{ptnb:eq:attn-Z-def}
    \tilde{v}_{k,d}^\kappa \;\coloneqq\; \sum_e [X_{k,e} W_{v,e,d}]^\kappa + [b_{v,d}]^\kappa,
    \qquad
    Z_{t,d}^\kappa \;\coloneqq\; \sum_{k} A_{t,k}\, \tilde{v}_{k,d}^\kappa
    \;=\; \sum_{k,e} A_{t,k}\, [X_{k,e} W_{v,e,d}]^\kappa + [b_{v,d}]^\kappa.
\end{align}
Note that $\tilde{v}_{k,d}^\kappa$ is the $\kappa$-stream split inherited from $W_v$ and the sign
of $X$ (not the standard scalar positive part $[v_{k,d}]^+ = \max(0, v_{k,d})$); the two coincide
only when $v_{k,d}$ has a single-signed predecessor decomposition.

\paragraph{Target rule.} Let $R^O_{t,d}$ be relevance arriving at $O_{t,d}$ from the surrounding LRP recursion. Apply the $\alpha\beta$-\ac{lrp} \ch{rule~\eqref{ptnb:eq:lrp-ab}} to the composite map $X \mapsto O$ in two stages, retaining $W_v$'s stream split $\tilde v_{k,d}^\kappa$ through~\eqref{ptnb:eq:attn-composite-weight}. \ch{Sequential $\alpha\beta$-\ac{lrp} on $X \to v \to O$ instead splits netted Values by sign and applies $(\alpha,\beta)$ twice, giving a different rule for mixed-sign predecessors. All reported ViT results and the implementation (Section~\ref{ptnb:app:vit-implementation}) use the composite rule; Section~\ref{ptnb:app:code-tests}'s tests check it against~\eqref{ptnb:eq:attn-target-input}.}

\paragraph{Zero-mass convention.}
All ratios of non-negative contributions in this subsection use
\begin{align}\label{ptnb:eq:attn-zero-ratio}
    \rho(u,D)\coloneqq
    \begin{cases}u/D,&D>0,\\0,&D=0,\end{cases}
    \qquad 0\leq u\leq D.
\end{align}
A zero denominator therefore gives zero. Zero-mass states route probability one to the cemetery and zero to ordinary successors, preserving a probability distribution. This applies to both $Z_{t,d}^\kappa$ and $\tilde v_{k,d}^\kappa$: zero $Z$ forces every $A_{t,k}\tilde v_{k,d}^\kappa$ to vanish; zero $\tilde v$ forces all matching input and bias contributions to vanish. The ideal formulas use $\varepsilon=0$. The released implementation also supports $Z_{t,d}^\kappa+\varepsilon$ with positive stabilization, routing the remaining share to the outside option.

\emph{AV mixing $O \to v$.} Applied to $O_{t,d} = \sum_k A_{t,k}\,v_{k,d}$, the rule gives the
relevance at the Value entry $v_{k,d}$ as the weighted difference of two per-stream contributions:
\begin{align}\label{ptnb:eq:attn-target-rule}
    R^{\mathrm{val}}_{k,d} \;=\; \alpha\,R^{\mathrm{val},+}_{k,d} \;-\;
    \beta\,R^{\mathrm{val},-}_{k,d},
    \qquad
    R^{\mathrm{val},\kappa}_{k,d} \;\coloneqq\; \sum_t \frac{A_{t,k}\,\tilde
    v_{k,d}^\kappa}{Z_{t,d}^\kappa}\,R^O_{t,d},
    \quad \kappa \in \{+,-\}.
\end{align}

\emph{Value projection $v \to X$.} Applied to $v_{k,d} = \sum_e W_{v,e,d}\,X_{k,e} + b_{v,d}$, the
rule routes each $\kappa$-stream contribution $R^{\mathrm{val},\kappa}_{k,d}$ to the input dim $e$
via the matching $\kappa$-stream ratio $[X_{k,e}W_{v,e,d}]^\kappa/\tilde v_{k,d}^\kappa$ (the bias
share is absorbed):
\begin{align}\label{ptnb:eq:attn-target-input-via-V}
    R^X_{k,e} \;=\; \sum_d \left(
    \alpha\, \frac{[X_{k,e} W_{v,e,d}]^+}{\tilde v_{k,d}^+}\,R^{\mathrm{val},+}_{k,d}
    \;-\;
    \beta\, \frac{[X_{k,e} W_{v,e,d}]^-}{\tilde v_{k,d}^-}\,R^{\mathrm{val},-}_{k,d}
    \right).
\end{align}

\emph{Aggregate.} Substitute~\eqref{ptnb:eq:attn-target-rule} into~\eqref{ptnb:eq:attn-target-input-via-V}. Positive $\tilde v_{k,d}^\kappa$ cancels; if it or $Z_{t,d}^\kappa$ is zero, both forms vanish by convention. Thus
\begin{align}\label{ptnb:eq:attn-target-input}
    R^X_{k,e}
    \;=\; \sum_{t,d} \left(
    \alpha\, \frac{A_{t,k}\, [X_{k,e} W_{v,e,d}]^+}{Z_{t,d}^+}
    \;-\;
    \beta\, \frac{A_{t,k}\, [X_{k,e} W_{v,e,d}]^-}{Z_{t,d}^-}
    \right) R^O_{t,d},
\end{align}
is exactly $\alpha\beta$-\ac{lrp} on the conditioned map $X \mapsto O$ with $W^{\mathrm{eff}}_{(k,e),(t,d)} = A_{t,k}\,W_{v,e,d}$ and bias $b_{v,d}$.

The state-space construction and the proof of Theorem~\ref{ptnb:thm:attn-ab-recovery} are given in Chapter~\ref{app:supp_attention} of the companion technical supplement.

\subsection{The reference as the policy of a QK oracle player}\label{ptnb:app:attention-oracle}

\S\ref{ptnb:app:attention-kl-routing} treated the softmax row $A_{t,\cdot}$ as a fixed reference policy.
It is itself the entropy-regularized best response of an auxiliary \emph{oracle} player whose
payoff is the detached query--key logit $e_{tk} = Q_t K_k^\top/\sqrt{d_h}$: by~\eqref{ptnb:eq:gibbs-lse}
at temperature $1$,
\begin{align}\label{ptnb:eq:attn-qk-oracle}
    \begin{gathered}
    A_{t,\cdot} \;=\; \argmax_{\pi \in \Delta(\{1,\dots,S\})} \Bigl\{ \textstyle\sum_{k} \pi(k)\,
    e_{tk} + \mathcal{H}(\pi) \Bigr\},\\
    \max_{\pi} \Bigl\{ \textstyle\sum_{k} \pi(k)\, e_{tk} + \mathcal{H}(\pi) \Bigr\}
    = \log \textstyle\sum_{k} \exp(e_{tk}).
    \end{gathered}
\end{align}
The oracle acts in the forward pass only and hands $\bar\pi_t = A_{t,\cdot}$ to the Value-routing
player; its states carry no occupation mass, so no relevance flows through $Q$ or $K$
(\S\ref{ptnb:app:vit-implementation}).

\subsection{Implementation details}\label{ptnb:app:vit-implementation}\label{sec:vit_implementation}

Four operator conventions align the \ch{\ac{rg} forward and backward passes on ViT-B/16} with \S\ref{ptnb:app:attention-state-space}'s explanatory walk. They define a DAVE-inspired \emph{conditioned forward pass}~\parencite{Wrobel2024} with the LayerNorm mean additionally frozen. Routing conservation includes bias and stabilizer absorption.

\paragraph{LayerNorm detachment.}\label{ptnb:app:dave-conditioned}
Standard backpropagation through $\text{LN}(x) = w_{\mathrm{LN}} \odot (x - m(x))/s(x) + b_{\mathrm{LN}}$ traverses data-dependent $m(x),s(x)$, \ch{coupling features within each token without mixing tokens}, and can violate \ac{lrp} conservation~\parencite{Ali2022AttnLRP}. Computing $m,s$ from $c = x.\text{detach}()$ gives the pointwise affine map $x \mapsto w_{\mathrm{LN}} \odot (x - m(c))/s(c) + b_{\mathrm{LN}}$. Our rule freezes both mean and scale; CP-LRP and DAVE freeze only the scale.

\paragraph{GELU gate detachment.}
Implement GELU as $x \cdot \Phi(c)$, with $c = x.\text{detach}()$ and standard normal CDF $\Phi$. Backward propagation then scales pointwise by constant $\Phi(c)$, realizing the elementwise LRP identity rule in modified Gradient$\times$Input implementations \parencite[Appendix~D.3]{Arras2025CloseLook}.

\paragraph{Softmax detachment.}
Softmax outputs are detached. As in \S\ref{ptnb:app:attention-oracle}, the cached row $A_{t,\cdot}$ is an auxiliary entropy-regularized equilibrium policy supplied to the Value-routing player; explanatory trajectories never enter $Q$ or $K$.

\paragraph{QK gradient detachment.}\label{ptnb:app:attention-grad-detach}
Softmax and QK-matmul backward hooks return zero, restricting QK to computing the detached KL reference $A_{t,\cdot}$; all attention-block explanatory mass follows Values. The composite rule of Theorem~\ref{ptnb:thm:attn-ab-recovery} retains each Value entry's projection-stream masses $\tilde v^\kappa_{k,d}$. Attention routes stream $\kappa$ proportionally to $A_{t,k}\tilde v^\kappa_{k,d}$, then the Value projection routes to same-sign $[X_{k,e}W_{v,e,d}]^\kappa$ contributions. Discounts $(\alpha,\beta)$ therefore act once per block, as in~\eqref{ptnb:eq:attn-target-input}. \ch{The released code implements only this rule (\texttt{av\_rule='merged'}); no reported result uses \S\ref{ptnb:app:attention-target-rule}'s sequential variant. Section~\ref{ptnb:app:code-tests}'s identity tests compare it with an independent reference for~\eqref{ptnb:eq:attn-target-input}.}

\ch{Together these conventions implement a conditioned \ac{rg} for \S\ref{ptnb:app:attention-target-rule}'s Value-only attribution question, shared with CP-LRP~\parencite{Ali2022AttnLRP}; AttnLRP~\parencite{Achtibat2024AttnLRP} also propagates relevance through softmax and QK.}

\subsection{Equivariant Spatial Averaging}\label{ptnb:app:equivariant}

The $16{\times}16$ patch tokenization breaks translation equivariance. Following DAVE~\parencite{Wrobel2024}, \ac{rg}+Eq explains transformed inputs, maps explanations back, and averages the aligned maps (a Monte-Carlo Reynolds average),
\begin{align}\label{ptnb:eq:equivariant-average}
    \mathcal{M}_N(x) \;=\; \frac{1}{N}\sum_{n=1}^{N} T_n^{-1}\bigl[\mathcal{M}(T_n x)\bigr],
\end{align}
where $\mathcal{M}$ is the \ac{rg} attribution and $T_n$ are i.i.d.\ random rotations (uniform in $\pm 10^\circ$), translations (up to $\pm 5\%$ per image axis) and horizontal flips (probability $0.5$), with $N = 50$ and no input noise. The conditioned pass (\S\ref{ptnb:app:vit-implementation}) separates operator variation from the effective transformation, following DAVE's argument~\parencite{Wrobel2024}. These transforms do not form a group, so this optional smoothing has no exactness guarantee.
\par\medskip
